% =============================================================================
% Discussion
% =============================================================================
\section{Discussion}\label{sec:discussion}

\subsection{Connection to Kac's problem}

Kac \cite{Kac1966} asked whether the eigenvalues of the Laplacian uniquely
determine the shape of a membrane.  Our result is an engineering analogue:
the eigenfrequencies of a \emph{viscoelastic shell} determine its shape and
stiffness, provided the model captures the geometric anisotropy.  The sphere
model fails precisely because it discards this anisotropy by collapsing all
geometry into a single radius.

Gordon, Webb, and Wolpert \cite{Gordon1992} showed that isospectral
non-isometric domains exist for membranes.  Our oblate shell model avoids
this pathology because the stiffness terms (bending, membrane, prestress)
create mode-dependent coupling to the geometry that is absent in the
Laplacian.

\subsection{Practical versus structural identifiability}

Following Raue et al.\ \cite{Raue2009}, we emphasise that our result is one
of \emph{practical identifiability}: given finite-precision measurements, the
parameter estimation is well-conditioned.  We do not claim global uniqueness
(structural identifiability), which would require proving injectivity of the
forward map over the entire parameter domain --- a significantly harder
problem that we leave for future work.

The condition number $\kappa \approx 83.5$ means that $1\%$ measurement noise
translates to at most $\sim 83.5\%$ worst-case parameter error in the least
observable direction.  In practice, the Cram\'er--Rao bounds show that the
\emph{expected} uncertainty is much smaller, of order a few percent for each
parameter.

\subsection{Implications for tissue characterisation}

The result suggests that vibroacoustic spectroscopy could, in principle, be
used for non-invasive assessment of organ stiffness and geometry.  The key
requirement is that the forward model must capture the organ's aspherical
shape.  An equivalent-sphere assumption, while computationally convenient,
destroys the information content needed for inversion.

This has particular relevance for:
\begin{itemize}
  \item \textbf{Liver fibrosis staging}, where organ stiffness increases
    with disease progression and the liver is markedly non-spherical.
  \item \textbf{Bladder volume estimation}, where the organ changes shape
    as it fills.
  \item \textbf{Non-destructive food testing}, as demonstrated by our
    watermelon validation case.
\end{itemize}

\subsection{Limitations}

Several limitations should be noted:
\begin{enumerate}
  \item We consider only three parameters $(a, c, E)$.  Including additional
    unknowns (e.g.\ wall thickness~$h$, Poisson's ratio~$\nu$, fluid
    density~$\rho_f$) would require more measured frequencies and may
    degrade conditioning.
  \item The forward model assumes a homogeneous, isotropic shell.  Real
    organs exhibit spatial heterogeneity and anisotropy.
  \item We assume the mode numbers of measured frequencies are known
    \emph{a priori}.  Mode identification from experimental data is itself
    a non-trivial problem.
  \item The noise model is Gaussian and frequency-proportional.  Real
    measurement noise may have different structure.
\end{enumerate}

\subsection{Future directions}

Natural extensions include:
\begin{itemize}
  \item Bayesian inversion with informative priors on tissue properties.
  \item Experimental validation using phantom organs with known properties.
  \item Extension to triaxial ellipsoids (three distinct semi-axes).
  \item Sensitivity analysis for the full parameter set using Sobol indices.
\end{itemize}
